\section{แนวทางพัฒนา ทดสอบ และนำขึ้นใช้งาน}
\subsection{แนวทางเพิ่ม Feature}
Backend ควรพัฒนาตามลำดับ DTO, service, controller, module และ schema/migration เมื่อจำเป็น แล้วทดสอบ validation, authentication, authorization และกรณีผิดพลาด Frontend ควรเพิ่ม type, API function หรือ hook, component/page และสถานะ loading, error, empty พร้อมตรวจ role และ responsive layout

เมื่อแก้ approval route ต้องรักษา snapshot ของคำร้องเก่า คำร้องที่สร้างแล้วไม่ควรเปลี่ยนเส้นทางย้อนหลังเพราะ master data ปัจจุบันเปลี่ยน

\subsection{API สำคัญ}
ทุก path อยู่ภายใต้ \code{API\_BASE\_PATH}
\begin{table}[H]\centering
  \begin{tabularx}{\textwidth}{|>{\ttfamily\fontsize{13pt}{16pt}\selectfont\raggedright\arraybackslash}p{2cm}|>{\ttfamily\fontsize{12pt}{15pt}\selectfont\raggedright\arraybackslash}p{8cm}|X|}\hline
    \tableheader{Method} & \tableheader{Path} & \tableheader{หน้าที่} \\\hline
    POST & /account/login-profile & อ่านหรือสร้าง profile หลัง login \\\hline
GET & /student/by-nontri/:id & อ่านข้อมูลนิสิต \\\hline
PATCH & /student/:id/contact & แก้ข้อมูลติดต่อ \\\hline
POST & /leave-request & สร้างคำร้อง \\\hline
POST & /leave-request/upload/evidence/:id & อัปโหลดหลักฐาน \\\hline
GET & /leave-request/:id & อ่านรายละเอียด \\\hline
GET & /leave-request/:id/approval-flow & อ่านขั้นอนุมัติ \\\hline
  \end{tabularx}
  \caption{API ในกระบวนการหลัก}
\end{table}

\begin{table}[H]\centering
  \begin{tabularx}{\textwidth}{|>{\ttfamily\fontsize{13pt}{16pt}\selectfont\raggedright\arraybackslash}p{2cm}|>{\ttfamily\fontsize{12pt}{15pt}\selectfont\raggedright\arraybackslash}p{8cm}|X|}\hline
    \tableheader{Method} & \tableheader{Path} & \tableheader{หน้าที่} \\\hline
    POST & /leave-request/:id/review & อนุมัติหรือไม่อนุมัติ \\\hline
POST & /leave-request/:id/lecturer-read & บันทึกการอ่านของอาจารย์ \\\hline
GET & /leave-request/:id/document & อ่านข้อมูลเอกสาร \\\hline
GET & /leave-request/public/:token & อ่านข้อมูลผ่าน signed token \\\hline
GET, PUT & /academic-setting/* & จัดการภาคการศึกษา \\\hline
GET, PUT & /auto-approve-setting/rules & จัดการกฎอนุมัติอัตโนมัติ \\\hline
GET & /warehouse-dashboard/live/:report & อ่านรายงานปัจจุบัน \\\hline
  \end{tabularx}
  \caption{API ในกระบวนการหลัก (ต่อ)}
\end{table}

ใช้ Postman collection ต่อไปนี้เป็นฐานทดสอบ
\begin{center}
  \code{apps/api/postman/}\quad
  \code{kukps-eng-leave-request-api.postman\_collection.json}
\end{center}
จากนั้นตั้ง base URL ให้รวม prefix ตรงกับ runtime

\subsection{ตรวจคุณภาพและ Build}
\begin{lstlisting}
pnpm lint
pnpm build
pnpm build:web
pnpm build:api
pnpm --filter api test:security
pnpm --filter api warehouse:etl:test
docker build -f Dockerfile.web -t kukps-eng-leave-request-web .
docker build -f Dockerfile.api -t kukps-eng-leave-request-api .
\end{lstlisting}
ควรทดสอบ login และ route ของทุก role, กระบวนการ create--approve/reject--lecturer read และไฟล์ JPG, JPEG, PNG, WEBP, PDF ทั้งกรณีผ่านและเกิน 2 MB ระบบไม่มี Docker Compose จึง deploy web, API และฐานข้อมูลแยกกัน

API จำกัด origin ด้วย \code{CORS\_ORIGIN} ซึ่งรับหลาย origin แบบคั่นด้วย comma ใน production ต้องระบุเฉพาะ web origin ที่อนุญาต ห้ามใช้ wildcard และต้องตั้ง \code{ROLE\_TEST=false}
